\hypertarget{appendix-b0efeb3f3648}{}
\section{Clustered lighting depth slicing}\label{appendix-b0efeb3f3648-clustered-lighting-depth-slicing}

Clustered lighting partitions the camera frustum into three-dimensional cells and assigns lights to the cells they can influence. A fragment then tests only the lights attached to its cell rather than every light in the scene.

The screen dimensions are naturally tiled in x and y. Depth needs more care because perspective scenes can span orders of magnitude. Uniform linear z slices often waste many clusters far from the camera or give poor resolution near it.

A common logarithmic mapping for positive camera-space depth \(z\) between near \(n\) and far \(f\) is

\[
t=
\frac{\log(z/n)}{\log(f/n)},
\]

then

\[
slice=
\left\lfloor tN_z\right\rfloor,
\]

clamped to \([0,N_z-1]\).

Equal multiplicative depth ratios occupy equal widths in log space. With \(n=0.1\) and \(f=1000\), the decades 1→10 and 10→100 span the same number of ideal log slices.

\subsection{Light assignment}\label{appendix-b0efeb3f3648-light-assignment}

Depth slicing alone does not make a clustered renderer. For each point/spot light, determine a conservative cluster range from its influence volume. Then append the light ID to every overlapping cluster.

The GPU implementation usually has two allocation problems:

\begin{enumerate}
\def\labelenumi{\arabic{enumi}.}
\tightlist
\item
  count how many light references each cluster needs;
\item
  prefix-scan those counts to obtain compact non-overlapping list offsets;
\item
  fill the light-index array.
\end{enumerate}

Atomic append alternatives exist, but they require overflow policy and contention measurement.

\subsection{False positives versus false negatives}\label{appendix-b0efeb3f3648-false-positives-versus-false-negatives}

A conservative cluster test may include a light that ultimately contributes nothing to a fragment. That costs shading work but preserves correctness. A false negative drops illumination and is a correctness bug.

Cluster bounds therefore need to match the exact camera projection/depth convention, including reversed-Z if hardware depth is used to derive z.

\subsection{Executable reference}\label{appendix-b0efeb3f3648-executable-reference}

tools/labs/realtime\_rendering\_reference\_lab.cpp validates the logarithmic slice mapping, near/far clamping, monotonicity and equal-decade behavior. It does not yet implement full 3D light-volume/cluster intersection or GPU compaction.

Related: \href{../01-Foundations-I/main.pdf\#nameddest=appendix-f2ee4d086415}{Prefix scan and stream compaction (01-Foundations-I, appendix-f2ee4d086415)}, \href{../03-Visual-Computing-I/main.pdf\#nameddest=appendix-275ec8847057}{Reversed-Z depth contracts (03-Visual-Computing-I, appendix-275ec8847057)}.
